\begin{table*}[t]
\centering
\scriptsize
\renewcommand{\arraystretch}{0.92}
\begin{tabularx}{\linewidth}{>{\raggedright\arraybackslash}p{0.17\linewidth}>{\raggedright\arraybackslash}p{0.27\linewidth}>{\raggedright\arraybackslash}p{0.25\linewidth}X}
\toprule
\textbf{Axis} & \textbf{Typical variation across papers} & \textbf{Why conclusions fail to transfer} & \textbf{Minimum reporting requirement} \\
\midrule
Retriever/generator setting & BM25 vs dense vs hybrid retrieval; small vs frontier generators; frozen vs fine-tuned models & A defense may work by compensating for one retriever or one model scale rather than solving the failure mode. & Retriever, index, top-$k$, reranker, generator, prompt, decoding, and context length. \\
Attack budget and access & Poisoned-document count, corpus write access, query control, retriever visibility, generator access & Attack success rates are incomparable when one attacker edits a corpus and another only controls the query. & Capability matrix, insertion budget, target set, stealth objective, and whether malicious evidence is retrieved. \\
Metric target & EM/F1, faithfulness, citation support, attack success, calibration error, selective risk, cost & Improving answer accuracy can worsen citation faithfulness; improving abstention can reduce coverage; reducing attacks can lower utility. & Report utility, evidence support, risk, calibration, cost, and security as a Pareto surface. \\
Statistical and reproducibility evidence & Single seed/model vs multiple seeds/models; code/data availability varies & Small prompt or model changes can reverse robustness rankings, especially for LLM-as-judge or attack benchmarks. & Seeds, confidence intervals or bootstrap tests where feasible, released prompts, data splits, and code or exact API settings. \\
Long-context baseline & RAG-only comparisons vs long-context LLMs with the same documents & A robustness gain may reflect retrieval compression rather than evidence reliability control. & Include a long-context or full-context baseline when context length permits, and report latency/cost. \\
Attribution vs answer accuracy & Some tasks score only final answer; others require evidence-supported citations & A model can answer correctly for the wrong reason, cite irrelevant evidence, or ignore retrieved support. & Claim-level support labels, citation precision/recall, and answer correctness reported separately. \\
\bottomrule
\end{tabularx}
\caption{Evaluation-assumption matrix for robust RAG. The table summarises why benchmark results often cannot be compared without aligning system setup, threat model, metric target, reproducibility evidence, and long-context baselines.}
\label{tab:evaluation-assumption-matrix}
\end{table*}

